\addcontentsline{toc}{section}{СПИСОК ИСПОЛЬЗОВАННЫХ ИСТОЧНИКОВ}

\begin{thebibliography}{9}
	
	\bibitem{1} Шакирьянов, Э. Д. Компьютерное зрение на Python. Первые шаги / Э. Д. Шакирьянов. -- Москва : Лаборатория знаний, 2025. -- 163 с. -- ISBN 978-5-00101-944-2. -- Текст : непосредственный.
	\bibitem{2} Сацюк, А. В. Компьютерное зрение и нейронные сети. Практика / А. В. Сацюк. -- Вологда : Инфра-Инженерия, 2025. -- 364 с. -- ISBN 978-5-9729-2706-7. -- Текст : непосредственный.
	\bibitem{3} Сацюк, А. В. Компьютерное зрение. Практика / А. В. Сацюк. -- Вологда : Инфра-Инженерия, 2025. -- 272 с. -- ISBN 978-5-9729-2346-5. -- Текст : непосредственный.
	\bibitem{4} Цехановский, В. В., Бутырский, Е. Ю., Жукова, Н. А., Мельников, В. Б. Искусственные нейронные сети. Практикум / В. В. Цехановский, Е. Ю. Бутырский, Н. А. Жукова, В. Б. Мельников. -- Москва : КноРус, 2026. -- 381 с. -- ISBN 978-5-406-15272-0. -- Текст : непосредственный.
	\bibitem{5} Ростовцев, В. С. Искусственные нейронные сети : учебник для вузов / В. С. Ростовцев. -- Санкт-Петербург : Лань, 2025. -- 216 с. -- ISBN 978-5-507-50568-5. -- Текст : непосредственный. 
	\bibitem{6} Николенко, С. И. Глубокое обучение: погружение в мир нейронных сетей / С. И. Николенко, А. Кадурин, Е. Архангельская. -- Санкт-Петербург : Питер, 2025. -- 476 с. -- ISBN 9785446115372. -- Текст : непосредственный.
	\bibitem{7} Гуриков, С. Р. Программирование на языке Python / С. Р. Гуриков. -- Москва : НИЦ ИНФРА-М, 2026. -- 332 с. -- ISBN 978-5-16-021562-4. -- Текст : непосредственный.
	\bibitem{8} Чернышев, С. А. Основы программирования на Python / С. А. Чернышев. -- Москва : Издательство Юрайт, 2026. -- 349 с. -- ISBN 978-5-534-17139-6. -- Текст : непосредственный.
	\bibitem{9} Федоров, Д. Ю. Программирование на Python / Д. Ю. Федоров. -- Москва : Издательство Юрайт, 2026. -- 166 с. -- ISBN 978-5-534-22181-7. -- Текст : непосредственный.
	\bibitem{10} Бильфельд, Н. В. Основы Python для будущих инженеров / Н. В. Бильфельд, А. В. Затонский. -- Вологда : Инфра-Инженерия, 2026. -- 468 с. -- ISBN 978-5-9729-3137-8. -- Текст : непосредственный.
	\bibitem{11} Никитина, Т. П. Программирование. Основы Python для инженеров / Т. П. Никитина, Л. В. Королев. -- Санкт-Петербург : Лань, 2026. -- 156 с. -- ISBN 978-5-507-51280-5. -- Текст : непосредственный.
	\bibitem{12} Дронов, В. А. Python 3 и PyQt 5. Разработка приложений / В. А. Дронов, Н. А. Прохоренок. -- Санкт-Петербург : БХВ-Петербург, 2018. -- 832 с. -- ISBN 978-5-9775-3648-6. -- Текст : непосредственный.
	\bibitem{13} Машнин, Т. Создание настольных Python приложений с графическим интерфейсом пользователя / Т. Машнин. -- [б. м.] : [б. и.], 2021. -- 147 с. -- ISBN 978-5-532-96908-7. -- Текст : непосредственный.
	\bibitem{14} Чаплыгин, А. А. Основы дипломного проектирования: учебное пособие для студентов направления подготовки 09.03.04 Программная инженерия, направленность (профиль) «Разработка программно-информационных систем» / А. А. Чаплыгин, В. В. Апальков, А. В. Малышев. -- Курск : Юго-Зап. гос. ун-т, 2022. -- 102 с. -- ISBN 978-5-907407-65-7. -- Текст : непосредственный.
	\bibitem{15} Кальб, И. Объектно-ориентированное программирование с помощью Python / И. Кальб. -- Москва : Эксмо, 2024. -- 512 с. -- ISBN 978-5-04-186627-3. -- Текст : непосредственный.
	\bibitem{16} Уилсон, Г. Разработка программного обеспечения на примерах с помощью Python / Г. Уилсон. -- Москва : Бомбора, 2025. -- 402 с. -- ISBN 978-5-04-237641-2. -- Текст : непосредственный.
	\bibitem{17} Мэтиз, Э. Изучаем Python: программирование игр, визуализация данных, веб-приложения / Э. Мэтиз. -- Санкт-Петербург : Питер, 2020. -- 512 с. -- ISBN 978-5-4461-1528-0. -- Текст : непосредственный.
	\bibitem{18} Левашов, П. Python с нуля / П. Левашов. -- СПб. : Питер, 2024. -- 448 с. -- ISBN 978-5-4461-2145-8. -- Текст : непосредственный.
	\bibitem{19} Прохоренок, Н. А. Python 3. Самое необходимое / Н. А. Прохоренок, В. А. Дронов. -- Санкт-Петербург : БХВ-Петербург, 2019. -- 608 с. -- ISBN 978-5-9775-3994-4. -- Текст : непосредственный.
	\bibitem{20} Белик, А. Г. Проектирование и архитектура программных систем : учебное пособие / А. Г. Белик, В. Н. Цыганенко. -- Омск : ОмГТУ, 2016. -- 96 с. -- ISBN 978-5-8149-2258-8. -- Текст : непосредственный.
	\bibitem{21} Вигерс, К. Разработка требований к программному обеспечению: практические приёмы сбора требований и управления ими при разработке программных продуктов / К. Вигерс, Д. Битти. -- Санкт-Петербург : БХВ, 2020. -- 736 с. -- ISBN 978-5-9775-3348-5. -- Текст : непосредственный.
	\bibitem{22} Юн, Ц. Рецепты Python. Коллекция лучших техник программирования / Ц. Юн. -- Санкт-Петербург : Питер, 2024. -- 512 с. -- ISBN 978-5-4461-2156-4. -- Текст : непосредственный.
	\bibitem{23} Бхаргава, А. Грокаем алгоритмы / А. Бхаргава. -- Санкт-Петербург : Питер, 2025. -- 210 с. -- ISBN 978-5-4461-4172-2. -- Текст : непосредственный.
	\bibitem{24} Python Documentation : официальный сайт. -- URL: https://www.python.org/doc/ (дата обращения: 16.05.2026). -- Текст : электронный.
	\bibitem{25} Шапиро, Л. Г. Компьютерное зрение : учебник для вузов / Л. Г. Шапиро, Д. В. Стокман. -- Санкт-Петербург : БХВ-Петербург, 2020. -- 756 с. -- ISBN 978-5-9963-3003-4. -- Текст : непосредственный.
	\bibitem{26} Семакин, И. Г. Основы программирования : учебник / И. Г. Семакин. -- Москва : Издательский центр «Академия», 2018. -- 432 с. -- ISBN 5-7695-2672-6. -- Текст : непосредственный.
	\bibitem{27} Демиденко, А. YOLO в действии: Обнаружение объектов / А. Демиденко. -- [б. м.] : [б. и.], 2025. -- URL: https://www.litres.ru/book/artem-demidenko/yolo-v-deystvii-obnaruzhenie-obektov-71679058/ (дата обращения: 10.06.2026).
	\bibitem{28} Матвеев, А. И. Цифровая обработка изображений в OpenCV. Практикум. Учебное пособие для вузов / А. И. Матвеев. -- Санкт-Петербург : Лань, 2026. -- 104 с. -- ISBN 978-5-507-51334-5. -- Текст : непосредственный.
	\bibitem{29} Qingliang, Z. Qt 5 and OpenCV 4 Computer Vision Projects / Z. Qingliang. -- Birmingham : Packt Publishing, 2019. -- 350 с. -- ISBN 9781789531831. -- Текст : непосредственный.
	\bibitem{30} Howse, J. Learning OpenCV 4 Computer Vision with Python 3 / J. Howse. -- Birmingham : Packt Publishing, 2020. -- 374 с. -- ISBN 9781789530643. -- Текст : непосредственный.
	\bibitem{31} Willman, J. Beginning PyQt / J. Willman. -- New York : Apress, 2022. -- 568 с. -- ISBN 9781484279984. -- Текст : непосредственный.
	\bibitem{32} Мартин, Р. Чистая архитектура. Искусство разработки программного обеспечения / Р. Мартин, пер. с англ. А. Кисилева. -- Санкт-Петербург: Питер,	2018. -- 351 с. -- ISBN 978-5-4461-0772-8. -- Текст: непосредственный.
	\bibitem{33} Фаулер, М. UML. Основы / М. Фаулер, пер. с англ. А. Петухова. -- 3-е изд. -- Санкт-Петербург : Символ-Плюс, 2004. -- 192 с. ISBN 978-5-93286-060-1. -- Текст : непосредственный.
	\bibitem{34} Шлее, М. Qt 5.10. Профессиональное программирование на C++ / М. Шлее. -- Санкт-Петербург : БХВ-Петербург, 2018. -- 1073 с. -- ISBN 978-5-9775-3678-3. -- Текст : непосредственный.
	\bibitem{35} Букунова, О. В. Разработка приложений с графическим пользовательским интерфейсом на языке Python. Учебное пособие для СПО / О. В. Букунова, С. В. Букунов. -- 3-е изд., стер. -- Санкт-Петербург : Лань, 2026. -- 89 с. -- ISBN 978-5-507-54799-9. -- Текст : непосредственный.
	\bibitem{36} Qt for Python : официальный сайт. -- URL: https://doc.qt.io/qtforpython-6/ (дата обращения: 16.04.2026). -- Текст : электронный.
	\bibitem{37} Ultralytics YOLO : официальный сайт. -- URL: https://docs.ultralytics.com (дата обращения: 10.04.2026). -- Текст : электронный.
	\bibitem{38} Куликова, И. В. Нейросети на Python. Основы ИИ и машинного обучения / И. В. Куликова. -- Санкт-Петербург : Наука и техника, 2025. -- 304 с. -- ISBN 978-5-907592-55-1. -- Текст : непосредственный.
	\bibitem{39} Гудфеллоу, Я. Глубокое обучение / Я. Гудфеллоу, И. Бенджио, А. Курвилль, пер. с англ. А. А. Слинкина. -- 2-е изд. -- Москва : ДМК Пресс, 2018. -- 652 с. -- ISBN 978-5-97060-618-6. -- Текст : непосредственный.
	\bibitem{40} Шолле, Ф. Глубокое обучение на Python / Ф. Шолле, пер. с англ. А. Киселева. -- 2-е межд. изд. -- Санкт-Петербург : Питер, 2023. -- 574 с. -- ISBN 978-5-4461-0770-4. -- Текст : непосредственный.
	\bibitem{41} Аггарвал, Ч. Нейронные сети и глубокое обучение : учебный курс / Ч. Аггарвал. -- Санкт-Петербург : Диалектика, 2020. -- 754 с. -- ISBN 9785907203013. -- Текст : непосредственный.
	\bibitem{42} Рашка, С. Python и машинное обучение. Машинное и глубокое обучение с использованием Python, scikit-learn и TensorFlow / С. Рашка, В. Мирджалили. -- Санкт-Петербург : Диалектика, 2021. -- 848 с. -- ISBN 978-5-907203-57-0. -- Текст : непосредственный.
	\bibitem{43} Солем, Я. Э. Программирование компьютерного зрения на Python / Я. Э. Солем, пер. с англ. А. А. Слинкина. -- Москва : ДМК Пресс, 2016. -- 312 с. -- ISBN 978-5-97060-200-3. -- Текст : непосредственный.
	\bibitem{44} Коул, А. Искусственный интеллект и компьютерное зрение. Реальные проекты на Python, Keras и TensorFlow / А. Коул, С. Ганджу, М. Казам, пер. с англ. А. Киселева. -- Санкт-Петербург : Питер, 2023. -- 624 с. -- ISBN 978-5-4461-1840-3. -- Текст : непосредственный.
	\bibitem{45} Клетте, Р. Компьютерное зрение. Теория и алгоритмы / Р. Клетте, пер. с англ. А. А. Слинкина. -- Москва : ДМК Пресс, 2019. -- 506 с. -- ISBN 978-5-97060-702-2. -- Текст : непосредственный.
	\bibitem{46} Корк, П. Машинное зрение и управление в робототехнике. Фундаментальные алгоритмы на Python / П. Корк, пер. с англ. А. Киселева. -- Москва : ДМК Пресс, 2026. -- 1106 с. -- ISBN 978-5-93700-429-1. -- Текст : непосредственный.
	\bibitem{47} Золкин, А. Л. Методы идентификации и отслеживания объектов в системах дополненной реальности и компьютерного зрения. Учебное пособие для вузов / А. Л. Золкин. -- Санкт-Петербург : Лань, 2025. -- 172 с. -- ISBN 978-5-507-53151-6. -- Текст : непосредственный.
	\bibitem{48} Селянкин, В. В. Компьютерное зрение. Анализ и обработка изображений. Учебное пособие для вузов / В. В. Селянкин. -- 4-е изд., стер. -- Санкт-Петербург : Лань, 2026. -- 152 с. -- ISBN 978-5-507-51201-0. -- Текст : непосредственный.
	\bibitem{49} Fitzpatrick, M. Create GUI Applications with Python \& Qt / M. Fitzpatrick. -- [б. м.] : [б. и.], 2025. -- 722 с. -- ISBN 9798287562120. -- Текст : непосредственный. 
	\bibitem{50} Jaganmohan, G. PySide GUI Application Development / G. Jaganmohan. -- Birmingham : Packt Publishing, 2016. -- 239 с. -- ISBN 9781785280368. -- Текст : непосредственный. 
	\bibitem{51} Chandrakar, S. Python GUI with PyQt / S. Chandrakar -- New Delhi : BPB Publications, 2023. -- 440 с. -- ISBN 9789355515575. -- Текст : непосредственный.
	
\end{thebibliography}